\documentclass[10pt,a4paper]{amsart}
\usepackage{amsmath,amssymb,amsthm,graphicx}
\usepackage{verbatim}
\newenvironment{reference}{\comment}{\endcomment}
%\newenvironment{reference}{}{}
\newenvironment{slogan}{\comment}{\endcomment}
\newenvironment{history}{\comment}{\endcomment}

% For commutative diagrams we use Xy-pic
\usepackage[all]{xy}

% We use 2cell for 2-commutative diagrams.
\xyoption{2cell}
% Theorem environments.
%
\theoremstyle{plain}
\newtheorem{theorem}[subsection]{Theorem}
\newtheorem{proposition}[subsection]{Proposition}
\newtheorem{lemma}[subsection]{Lemma}

\theoremstyle{definition}
\newtheorem{definition}[subsection]{Definition}
\newtheorem{example}[subsection]{Example}
\newtheorem{exercise}[subsection]{Exercise}
\newtheorem{situation}[subsection]{Situation}

\theoremstyle{remark}
\newtheorem{remark}[subsection]{Remark}
\newtheorem{remarks}[subsection]{Remarks}

\numberwithin{equation}{subsection}

% Macros
%
\def\lim{\mathop{\mathrm{lim}}\nolimits}
\def\colim{\mathop{\mathrm{colim}}\nolimits}
\def\Spec{\mathop{\mathrm{Spec}}}
\def\Hom{\mathop{\mathrm{Hom}}\nolimits}
\def\Ext{\mathop{\mathrm{Ext}}\nolimits}
\def\SheafHom{\mathop{\mathcal{H}\!\mathit{om}}\nolimits}
\def\SheafExt{\mathop{\mathcal{E}\!\mathit{xt}}\nolimits}
\def\Sch{\mathit{Sch}}
\def\Mor{\mathop{\mathrm{Mor}}\nolimits}
\def\Ob{\mathop{\mathrm{Ob}}\nolimits}
\def\Sh{\mathop{\mathit{Sh}}\nolimits}
\def\NL{\mathop{N\!L}\nolimits}
\def\CH{\mathop{\mathrm{CH}}\nolimits}
\def\proetale{{pro\text{-}\acute{e}tale}}
\def\etale{{\acute{e}tale}}
\def\QCoh{\mathit{QCoh}}
\def\Ker{\mathop{\mathrm{Ker}}}
\def\Im{\mathop{\mathrm{Im}}}
\def\Coker{\mathop{\mathrm{Coker}}}
\def\Coim{\mathop{\mathrm{Coim}}}

% Boxtimes
%
\DeclareMathSymbol{\boxtimes}{\mathbin}{AMSa}{"02}

%
% Macros for moduli stacks/spaces
%
\def\QCohstack{\mathcal{QC}\!\mathit{oh}}
\def\Cohstack{\mathcal{C}\!\mathit{oh}}
\def\Spacesstack{\mathcal{S}\!\mathit{paces}}
\def\Quotfunctor{\mathrm{Quot}}
\def\Hilbfunctor{\mathrm{Hilb}}
\def\Curvesstack{\mathcal{C}\!\mathit{urves}}
\def\Polarizedstack{\mathcal{P}\!\mathit{olarized}}
\def\Complexesstack{\mathcal{C}\!\mathit{omplexes}}
% \Pic is the operator that assigns to X its picard group, usage \Pic(X)
% \Picardstack_{X/B} denotes the Picard stack of X over B
% \Picardfunctor_{X/B} denotes the Picard functor of X over B
\def\Pic{\mathop{\mathrm{Pic}}\nolimits}
\def\Picardstack{\mathcal{P}\!\mathit{ic}}
\def\Picardfunctor{\mathrm{Pic}}
\def\Deformationcategory{\mathcal{D}\!\mathit{ef}}
\usepackage{cleveref}
\crefname{originallemma}{original lemma}{original lemmas}
\Crefname{originallemma}{Original lemma}{Original lemmas}
\makeatletter
\newlabel{cohomology-lemma-exact-sequence-j-star}{{08BV}{source}{}{sourcepointer-1}{}}
\newlabel{cohomology-lemma-exact-sequence-j-star@cref}{{[originallemma][0][]08BV}{[1][0][]source}{}{}{}}
\makeatother
\crefname{originallemma}{original lemma}{original lemmas}
\Crefname{originallemma}{Original lemma}{Original lemmas}
\makeatletter
\newlabel{derived-lemma-essentially-constant-2-out-of-3}{{0G3A}{source}{}{sourcepointer-2}{}}
\newlabel{derived-lemma-essentially-constant-2-out-of-3@cref}{{[originallemma][0][]0G3A}{[1][0][]source}{}{}{}}
\makeatother
\crefname{originallemma}{original lemma}{original lemmas}
\Crefname{originallemma}{Original lemma}{Original lemmas}
\makeatletter
\newlabel{cohomology-lemma-mayer-vietoris-hom}{{08BW}{source}{}{sourcepointer-3}{}}
\newlabel{cohomology-lemma-mayer-vietoris-hom@cref}{{[originallemma][0][]08BW}{[1][0][]source}{}{}{}}
\makeatother
\crefname{originallemma}{original lemma}{original lemmas}
\Crefname{originallemma}{Original lemma}{Original lemmas}
\makeatletter
\newlabel{lemma-lift-map}{{0G2H}{source}{}{sourcepointer-4}{}}
\newlabel{lemma-lift-map@cref}{{[originallemma][0][]0G2H}{[1][0][]source}{}{}{}}
\makeatother
\crefname{originallemma}{original lemma}{original lemmas}
\Crefname{originallemma}{Original lemma}{Original lemmas}
\makeatletter
\newlabel{lemma-extension-by-zero}{{0G4L}{source}{}{sourcepointer-5}{}}
\newlabel{lemma-extension-by-zero@cref}{{[originallemma][0][]0G4L}{[1][0][]source}{}{}{}}
\makeatother
\crefname{originalreference}{original reference}{original references}
\Crefname{originalreference}{Original reference}{Original references}
\makeatletter
\newlabel{categories-remark-pro-category-copresheaves}{{0G2X}{source}{}{sourcepointer-6}{}}
\newlabel{categories-remark-pro-category-copresheaves@cref}{{[originalreference][0][]0G2X}{[1][0][]source}{}{}{}}
\makeatother
\title{Being pro-isomorphic to a Deligne system is local on an open covering}
\author{The Stacks Project Authors}
\date{stacks-project:duality.tex:local-snapshot-20261012; GNU FDL 1.2 or later}
\begin{document}
\maketitle
\noindent\textit{Editorial note.} All delivered mathematical text is the author's own, copied from the pinned Stacks Project chapter file duality.tex; no statement or proof step was rewritten, reordered or removed.\par\medskip
\noindent\textit{Editorial note.} Preamble provenance: the pinned chapter file opens with an input of the shared preamble and carries no preamble text of its own, so the project's single Stacks-wide preamble staged at sources/stacks-duality/preamble.tex is excerpted verbatim under the recipe fixed by LANE\_BRIEF 22.1, as three spans, lines 8-13 (verbatim package and the reference/slogan/history environments), lines 14-18 (Xy-pic setup) and lines 154-217 (the native theorem declarations and mathematical macros). Nothing was added, removed or redefined; the wrapper class and the amsmath/amssymb/amsthm/graphicx base packages are supplied by the assembler.\par\medskip
\noindent\textit{Editorial note.} The counted statement and proof contain no cross reference, no citation and no included file, so the delivered body is self contained.\par\medskip
\noindent\textit{Editorial note.} Canonical location of the counted proposition: Stacks tag 0G4T (sources/stacks-duality/pinned-original-source.tex line 8358).\par\medskip
\noindent\textit{Allowed standard prerequisites (editorial).}\begin{itemize}
\item Deligne systems and their restrictions to opens. Same chapter, the definition of a Deligne system preceding this lemma, invoked in the statement and at proof lines 8384-8386. Fixes the property being tested locally.
\item The colimit of Hom out of a Deligne system computes Hom on the open. Same chapter, Lemma labelled lemma-lift-map, cited at proof lines 8408-8412. Makes all but the middle horizontal arrow of the diagram an isomorphism.
\item Mayer-Vietoris for Hom and essential constancy along distinguished triangles. Cohomology, Lemma labelled cohomology-lemma-mayer-vietoris-hom, cited at proof lines 8404-8406, and Cohomology, Lemma labelled cohomology-lemma-exact-sequence-j-star with Derived Categories, Lemma labelled derived-lemma-essentially-constant-2-out-of-3, cited at proof lines 8388-8392. Give the exactness of the vertical sequences and the essential constancy of the system on U.
\item Existence of a Deligne system with prescribed constant value, and pro-isomorphism through colimits of Hom. Same chapter, Lemma labelled lemma-extension-by-zero, cited at proof lines 8418-8420, and Categories, Remark labelled categories-remark-pro-category-copresheaves, cited at proof lines 8422-8425. Concludes that the given system is pro-isomorphic to a Deligne system.
\end{itemize}\medskip
\ifdefined\hypertarget\hypertarget{sourcepointer-1}{}\fi
\noindent\textit{Source reference.} Original lemma 08BV: sources/stacks-chow/tags.txt, source lines 10580--10580. Cross-chapter standard prerequisite located in the project tag list staged in the local snapshot; invoked at proof line 8388 of the counted proof. It is retained only as a navigation pointer; every load-bearing step is inside the delivered span.\par\medskip
\ifdefined\hypertarget\hypertarget{sourcepointer-2}{}\fi
\noindent\textit{Source reference.} Original lemma 0G3A: sources/stacks-chow/tags.txt, source lines 20085--20085. Cross-chapter standard prerequisite located in the project tag list staged in the local snapshot; invoked at proof line 8389 of the counted proof. It is retained only as a navigation pointer; every load-bearing step is inside the delivered span.\par\medskip
\ifdefined\hypertarget\hypertarget{sourcepointer-3}{}\fi
\noindent\textit{Source reference.} Original lemma 08BW: sources/stacks-chow/tags.txt, source lines 10581--10581. Cross-chapter standard prerequisite located in the project tag list staged in the local snapshot; invoked at proof line 8412 of the counted proof. It is retained only as a navigation pointer; every load-bearing step is inside the delivered span.\par\medskip
\ifdefined\hypertarget\hypertarget{sourcepointer-4}{}\fi
\noindent\textit{Source reference.} Original lemma 0G2H: sources/stacks-duality/pinned-original-source.tex, source lines 7920--7920. Same-chapter pointer into the author environment lemma at source line 7920; invoked at proof line 8415 of the counted proof. It is retained only as a navigation pointer; every load-bearing step is inside the delivered span.\par\medskip
\ifdefined\hypertarget\hypertarget{sourcepointer-5}{}\fi
\noindent\textit{Source reference.} Original lemma 0G4L: sources/stacks-duality/pinned-original-source.tex, source lines 8103--8103. Same-chapter pointer into the author environment lemma at source line 8103; invoked at proof line 8420 of the counted proof. It is retained only as a navigation pointer; every load-bearing step is inside the delivered span.\par\medskip
\ifdefined\hypertarget\hypertarget{sourcepointer-6}{}\fi
\noindent\textit{Source reference.} Original reference 0G2X: sources/stacks-chow/tags.txt, source lines 20072--20072. Cross-chapter standard prerequisite located in the project tag list staged in the local snapshot; invoked at proof line 8424 of the counted proof. It is retained only as a navigation pointer; every load-bearing step is inside the delivered span.\par\medskip
\begin{lemma}
\label{lemma-characterize-extension-by-zero}
Let $j : U \to X$ be an open immersion of Noetherian schemes. Let
$(K_n)$ be an inverse system in $D^b_{\textit{Coh}}(\mathcal{O}_X)$.
Let $X = W_1 \cup \ldots \cup W_r$ be an open covering.
The following are equivalent
\begin{enumerate}
\item $(K_n)$ is pro-isomorphic to a Deligne system,
\item for each $i$ the restriction $(K_n|_{W_i})$
is pro-isomorphic to a Deligne system with respect to
the open immersion $U \cap W_i \to W_i$.
\end{enumerate}
\end{lemma}

\begin{proof}
By induction on $r$. If $r = 1$ then the result is clear.
Assume $r > 1$. Set $V = W_1 \cup \ldots \cup W_{r - 1}$.
By induction we see that $(K_n|_V)$ is a Deligne system.
This reduces us to the discussion in the next paragraph.

\medskip\noindent
Assume $X = V \cup W$ is an open covering and
$(K_n|_W)$ and $(K_n|_V)$ are pro-isomorphic to Deligne systems.
We have to show that $(K_n)$ is pro-isomorphic to a Deligne system.
Observe that $(K_n|_{V \cap W})$ is pro-isomorphic to a Deligne system
(it follows immediately from the construction of Deligne systems
that restrictions to opens preserves them). In particular the pro-systems
$(K_n|_{U \cap V})$,
$(K_n|_{U \cap W})$, and
$(K_n|_{U \cap V \cap W})$
are essentially constant. It follows from the distinguished triangles
in Cohomology, Lemma \ref{cohomology-lemma-exact-sequence-j-star} and
Derived Categories, Lemma \ref{derived-lemma-essentially-constant-2-out-of-3}
that $(K_n|_U)$ is essentially constant.
Denote $K \in D^b_{\textit{Coh}}(\mathcal{O}_U)$ the value of this system.
Let $L$ be an object of $D^b_{\textit{Coh}}(\mathcal{O}_X)$.
Consider the diagram
$$
\xymatrix{
\colim \Ext^{-1}(K_n|_V, L|_V) \oplus
\colim \Ext^{-1}(K_n|_W, L|_W) \ar[r] \ar[d] &
\Ext^{-1}(K|_{U \cap V}, L|_{U \cap V}) \oplus
\Ext^{-1}(K|_{U \cap W}, L|_{U \cap W}) \ar[d] \\
\colim \Ext^{-1}(K_n|_{V \cap W}, L|_{V \cap W}) \ar[r] \ar[d] &
\Ext^{-1}(K|_{U \cap V \cap W}, L|_{U \cap V \cap W}) \ar[d] \\
\colim \Hom(K_n, L) \ar[d] \ar[r] &
\Hom(K|_U, L|_U) \ar[d] \\
\colim \Hom(K_n|_V, L|_V) \oplus \colim \Hom(K_n|_W, L|_W) \ar[r] \ar[d] &
\Hom(K|_{U \cap V}, L|_{U \cap V}) \oplus
\Hom(K|_{U \cap W}, L|_{U \cap W}) \ar[d] \\
\colim \Hom(K_n|_{V \cap W}, L|_{V \cap W}) \ar[r] &
\Hom(K|_{U \cap V \cap W}, L|_{U \cap V \cap W})
}
$$
The vertical sequences are exact by
Cohomology, Lemma \ref{cohomology-lemma-mayer-vietoris-hom}
and the fact that filtered colimits are exact.
All horizontal arrows except for the middle one are isomorphisms
by Lemma \ref{lemma-lift-map} and the fact that pro-isomorphic systems
have the same colimits. Hence the middle one is an isomorphism too by
the 5-lemma. It follows that $(K_n)$ is pro-isomorphic to
a Deligne system for $K$. Namey, if $(K'_n)$ is a Deligne system
whose restriction to $U$ has constant value $K$ (which
exists by Lemma \ref{lemma-extension-by-zero}), then
we have $\colim \Hom_X(K'_n, L) = \Hom_U(K, L|_U)$ as well.
Hence the pro-systems $(K_n)$ and $(K'_n)$ are
pro-isomorphic by
Categories, Remark \ref{categories-remark-pro-category-copresheaves}.
\end{proof}

\end{document}
